% arara: xelatex
\documentclass[12pt]{article}

% \usepackage{physics}

\usepackage{hyperref}
\hypersetup{
    colorlinks=true,
    linkcolor=blue,
    filecolor=magenta,      
    urlcolor=cyan,
    pdftitle={Overleaf Example},
    pdfpagemode=FullScreen,
    }

\usepackage{verse}

\usepackage{tikzducks}

\usepackage{tikz} % картинки в tikz
\usetikzlibrary{shapes, arrows, positioning}
\usepackage{microtype} % свешивание пунктуации

\usepackage{array} % для столбцов фиксированной ширины

\usepackage{indentfirst} % отступ в первом параграфе

\usepackage{sectsty} % для центрирования названий частей
\allsectionsfont{\centering}

\usepackage{amsmath, amsfonts, amssymb} % куча стандартных математических плюшек

\usepackage{comment}

\usepackage[top=2cm, left=1.2cm, right=1.2cm, bottom=2cm]{geometry} % размер текста на странице

\usepackage{lastpage} % чтобы узнать номер последней страницы

\usepackage{enumitem} % дополнительные плюшки для списков
%  например \begin{enumerate}[resume] позволяет продолжить нумерацию в новом списке
\usepackage{caption}

\usepackage{url} % to use \url{link to web}

\usepackage{framed} % для рамок и черты слева от минитеории, \leftbar

\newcommand{\smallduck}{\begin{tikzpicture}[scale=0.3]
    \duck[
        cape=black,
        hat=black,
        mask=black
    ]
    \end{tikzpicture}}

\usepackage{fancyhdr} % весёлые колонтитулы
\pagestyle{fancy}
\lhead{ШВН-2026}
\chead{}
\rhead{Теория невероятностей}
\lfoot{}
\cfoot{}
%\rfoot{\begin{tikzpicture}[scale=0.4]
%    \duck[hockey=brown!70!black, 
%    crazyhair=red]
%\end{tikzpicture}}

\rfoot{\shuffleducks
    \begin{tikzpicture}[scale=0.4] 
        \duck[umbrella=cyan, 
        speech={\thepage},
        bubblecolour=cyan!20!white,
        laughing
        ],  
    \end{tikzpicture}} 


\renewcommand{\headrulewidth}{0.4pt}
\renewcommand{\footrulewidth}{0.4pt}

\usepackage{tcolorbox} % рамочки!

\usepackage{todonotes} % для вставки в документ заметок о том, что осталось сделать
% \todo{Здесь надо коэффициенты исправить}
% \missingfigure{Здесь будет Последний день Помпеи}
% \listoftodos - печатает все поставленные \todo'шки


% более красивые таблицы
\usepackage{booktabs}
% заповеди из докупентации:
% 1. Не используйте вертикальные линни
% 2. Не используйте двойные линии
% 3. Единицы измерения - в шапку таблицы
% 4. Не сокращайте .1 вместо 0.1
% 5. Повторяющееся значение повторяйте, а не говорите "то же"


\setcounter{MaxMatrixCols}{20}
% by crazy default pmatrix supports only 10 cols :)


\usepackage{fontspec}
\usepackage{libertine}
\usepackage{polyglossia}

\setmainlanguage{russian}
\setotherlanguages{english}

% download "Linux Libertine" fonts:
% http://www.linuxlibertine.org/index.php?id=91&L=1
% \setmainfont{Linux Libertine O} % or Helvetica, Arial, Cambria
% why do we need \newfontfamily:
% http://tex.stackexchange.com/questions/91507/
% \newfontfamily{\cyrillicfonttt}{Linux Libertine O}

\AddEnumerateCounter{\asbuk}{\russian@alph}{щ} % для списков с русскими буквами
% \setlist[enumerate, 2]{label=\asbuk*),ref=\asbuk*}

%% эконометрические сокращения
\DeclareMathOperator{\Cov}{\mathbb{C}ov}
\DeclareMathOperator{\Corr}{\mathbb{C}orr}
\DeclareMathOperator{\Var}{\mathbb{V}ar}
\DeclareMathOperator{\col}{col}
\DeclareMathOperator{\row}{row}

\let\P\relax
\DeclareMathOperator{\P}{\mathbb{P}}

\DeclareMathOperator{\E}{\mathbb{E}}
% \DeclareMathOperator{\tr}{trace}
\DeclareMathOperator{\card}{card}
\DeclareMathOperator{\mgf}{mgf}

\DeclareMathOperator{\Convex}{Convex}
\DeclareMathOperator{\plim}{plim}

\usepackage{mathtools}
\DeclarePairedDelimiter{\norm}{\lVert}{\rVert}
\DeclarePairedDelimiter{\abs}{\lvert}{\rvert}
\DeclarePairedDelimiter{\scalp}{\langle}{\rangle}
\DeclarePairedDelimiter{\ceil}{\lceil}{\rceil}

\newcommand{\cN}{\mathcal{N}}
\newcommand{\cF}{\mathcal{F}}

\newcommand{\RR}{\mathbb{R}}
\newcommand{\NN}{\mathbb{N}}
\newcommand{\hb}{\hat{\beta}}
\newcommand{\dPois}{\mathrm{Pois}}



\newcommand{\addtag}[1]{\index{#1}}

\begin{document}


\begin{verse}
    \begin{flushright}
        — Что мы знаем о лисе? \\
        — Ничего. И то не все. \\

        Борис Заходер
    \end{flushright}
\end{verse}

\section{Встреча первая}

\begin{leftbar}
Константы обозначаем строчными английскими буквами: $a$, $x$, События ­— заглавными английскими буквами начала алфавита $A$, $B$, $C$, $D$, \ldots{ } Вероятность события — $\P(A)$.
Случайные величины обозначаем заглавными английскими буквы второй половины и конца алфавита: $X$, $Y$, $W$, $Z$, \ldots{ } Математическое ожидание (среднее значение) — $\E(X)$. 
\end{leftbar}

\begin{enumerate}
    \item Джон Конвей \raisebox{-3pt}{\begin{tikzpicture}[scale=0.3]
    \duck[
        cape=black,
        hat=black,
        mask=black
    ]
    \end{tikzpicture}} подбрасывает правильную монетку до выпадения трёх орлов подряд. 

    Обозначим с помощью $X$ число потребовавшихся бросков. 

    \begin{enumerate}
        \item Найдите вероятность $\P(X = 3)$, $\P(X = 4)$.
        \item Найдите математическое ожидание $\E(X)$.
    \end{enumerate}

\item Маша и Даша по очереди подбрасывают кубик\addtag{кубик}.
Посуду будет мыть тот, кто первым выбросит шестерку.
Маша бросает кубик первой.

\begin{enumerate}
\item Какова вероятность того, что посуду будет мыть Маша?
\item Сколько в среднем раз они будут бросать кубик\addtag{кубик}?
\item Сколько в среднем конфет съест Даша, если за каждую единичку, выброшенную Машей, Даша будет съедать одну конфету, а за двойку — две?
\end{enumerate}


\item Маша и Даша играют в следующую игру. 
 Правильный кубик подкидывают неограниченное число раз.
 Если на кубике\addtag{кубик} выпадает $1$, то казино добавляет в шляпу один рубль и игра продолжается.
 Если на кубике выпадает $2$, то казино увеличивает сумму в шляпе в два раза и игра продолжается. 
 Если на кубике выпадает $3$, то казино выплачивает Маше $3$ рубля и игра продолжается. 
 Если на кубике выпадает $4$, то казино выплачивает Даше $4$ рубля и игра продолжается. 
 Если на кубике выпадает $5$, то сумма в шляпе сгорает и игра продолжается.
 Если на кубике выпадает $6$, то Маша и Даша получают сумму в шляпе, делят её поровну и игра оканчивается.
Изначально на столе стоит пустая шляпа.
 \begin{enumerate}
\item Какова ожидаемая продолжительность игры?
\item Чему равен ожидаемый выигрыш Маши и ожидаемый выигрыш Даши?
\item Чему равны ожидаемые расходы организаторов игры?
\item Чему равен ожидаемый выигрыш Маши, если изначально в шляпе лежит $100$ рублей?
\end{enumerate}

\item Илье Муромцу\addtag{герои!Илья Муромец} предстоит дорога к камню. 
От камня начинаются ещё три дороги.
Каждая из тех дорог снова оканчивается камнем. 
И от каждого камня начинаются ещё три дороги.
И каждые те три дороги оканчиваются камнем\ldots{ } 
И так далее до бесконечности.
На каждой дороге живёт трёхголовый Змей Горыныч\addtag{герои!Змей Горыныч}.
Каждый Змей Горыныч бодрствует независимо от других с вероятностью (хм, Вы не поверите!) одна третья.
%У Василисы Премудрой\addtag{герои!Василиса Премудрая} существует Чудо-Карта, на которой видно,
% какие Змеи Горынычи бодрствуют\addtag{герои!Змей Горыныч}, а какие — нет.

\begin{enumerate}
\item Какова вероятность того, что Илья Муромец будет постоянно идти мимо спящих Горынычей, если будет выбирать каждый ход случайно?
\item Какова вероятность того,
что Василиса Премудрая\addtag{герои!Василиса Премудрая}, знающая состояние всех Горынычей, 
сможет указать Ильи Муромцу бесконечный жизненный путь проходящий исключительно мимо спящих Горынычей?
\end{enumerate}

\item Маша и Даша подкидывают одну правильную монетку. 
Маша выигрывает, как только выпадет последовательность $HTH$, Даша — как только выпадет последовательность $HHT$.
Игра заканчивается, как только выпадает одна из этих последовательностей. 

\begin{enumerate}
    \item Какова вероятность того, что игра длится вечно?
    \item Сколько в среднем длится одна игра?
    \item Какова вероятность того, что выиграет Маша?
\end{enumerate}

\end{enumerate}

\newpage
\section{Встреча вторая}

\begin{enumerate}
    \item Эдвард «Чёрная борода» \raisebox{-3pt}{\begin{tikzpicture}[scale=0.3]
    \duck[
        cape=black,
        hat=black,
        mask=black
    ]
    \end{tikzpicture}} подбрасывает правильную монетку до выпадения трёх орлов \emph{не обязательно подряд}. 

    Обозначим с помощью $Y$ число потребовавшихся бросков. 

    \begin{enumerate}
        \item Найдите вероятность $\P(Y = 3)$, $\P(Y = 4)$.
        \item Найдите математическое ожидание $\E(Y)$.
    \end{enumerate}

    \item Джек Воробей \raisebox{-3pt}{\begin{tikzpicture}[scale=0.3]
    \duck[
        cape=black,
        hat=black,
        mask=black
    ]
    \end{tikzpicture}} подбрасывает правильную монетку до выпадения трёх орлов подряд. 
    
    Однако на $3$-м, $6$-м, $9$-м и так далее бросках Джека отвлекает попугай, и Джек не может остановиться в эти моменты времени.
    Например, на последовательности $THTHHHH$ остановка произойдёт после $7$-го броска. 

    Сколько бросков в среднем пройдёт пока Джек заметит трёх выпавшех орлов?

    \item Джон Сильвер  \raisebox{-3pt}{\begin{tikzpicture}[scale=0.3]
    \duck[
        cape=black,
        hat=black,
        mask=black
    ]
    \end{tikzpicture}} приходит в таверну «Огненная зебра» и заказывает сразу три пинты эля. 
    Хитрый бармен «Огненной зебры» разбавляет каждую пинту эля с вероятностью $0.5$. 
    Если больше половины из выпитых пинт эля разбавлены, то Джон разносит таверну к чертям собачьим. 
    Если меньше, то Джон заказывает очередную пинту эля. 
    Из-за пагубного пристрастия к алкоголю у Джона проблемы с памятью и он помнит только последние три пинты. 

    Сколько в среднем эля выпьет Джон Сильвер пока не разнесёт таверну?

    \item Сэр Генри Морган \raisebox{-3pt}{\begin{tikzpicture}[scale=0.3]
    \duck[
        cape=black,
        hat=black,
        mask=black
    ]
    \end{tikzpicture}} подбрасывает монетку\addtag{монетка} 
 неограниченной количество раз до тех пор, пока число орлов не окажется равным числу решек (не считая начала игры).
    Монетка\addtag{монетка} выпадает орлом с вероятностью $p$.

    \begin{enumerate}
        \item Какова вероятность того, что сэр Генри Морган будет подбрасывать монетку вечно?
        \item Сколько в среднем бросков пройдёт до окончания игры?
    \end{enumerate}

    \item Бартоломью Робертс \raisebox{-3pt}{\begin{tikzpicture}[scale=0.3]
    \duck[
        cape=black,
        hat=black,
        mask=black
    ]
    \end{tikzpicture}} подбрасывает монетку\addtag{монетка} 
 неограниченной количество раз до тех пор, пока число орлов не окажется в два раза больше числа решек (не считая начала игры).
    Монетка\addtag{монетка} выпадает орлом с вероятностью $p$.

    \begin{enumerate}
        \item Какова вероятность того, что Бартоломью будет подбрасывать монетку вечно?
        \item Сколько в среднем бросков пройдёт до окончания игры?
    \end{enumerate}

\end{enumerate}

\newpage
\section{Встреча третья}
\begin{enumerate}
    \item Леонардо Пизанский выращивает кроликов. 
    Изначально у него одна юная пара кроликов $F_1$.
    Юная пара кроликов становится взрослой за один месяц. 
    Каждая взрослая пара кроликов приносит новую юную пару кроликов каждый месяц. 
    \begin{enumerate}
        \item Сколько пар будет у Леонардо в ближайшие месяцы? Найдите $F_2$, $F_3$, $F_4$, $F_5$.
        \item Нарисуйте граф для слов-родословных кроликов Леонардо. 
        \item Найдите общую формулу для $F_n$.
    \end{enumerate}
    
    \item Пуанкаре подбрасывает монетку до выпадения последовательности $HTH$.

\begin{enumerate}[label=\alph*)]
    \item Нарисуйте граф для этой игры. 
    \item Выпиши систему уравнений на слова-траектории от старта до заданной точки. 
    \item Выпиши систему уравнений на слова-траектории от заданной точки до финиша. 
    \item Выпиши систему уравнений на слова-траектории от заданной точки и недошедшие до финиша. 
    \item Объясните, как превратить слова-траетории в осмысленные вероятности или средние значения.
\end{enumerate}

\item Маша и Даша подкидывают одну правильную монетку. 
Маша выигрывает, как только выпадет последовательность $HHH$, Даша — как только выпадет последовательность $HTH$.
Игра заканчивается, как только выпадает одна из этих последовательностей. 

\begin{enumerate}[label=\alph*)]
    \item Нарисуйте граф для этой игры. 
    \item Выпиши систему уравнений на слова-траектории от старта до заданной точки. 
    \item Выпиши систему уравнений на слова-траектории от заданной точки до финиша в $HHH$. 
    \item Выпиши систему уравнений на слова-траектории от заданной точки до финиша в $HTH$. 
    \item Выпиши систему уравнений на слова-траектории от заданной точки и недошедшие до финиша. 
    \item Сколько в среднем времени длится игра?
    \item Какова вероятность того, что выиграет Маша?
\end{enumerate}


\end{enumerate}


\newpage
\section{Встреча четвёртая}
\begin{enumerate}
\item Эльчалан Моссель подкидывает кубик до первой шестёрки. 
\begin{enumerate}
    \item Сколько в среднем бросков ему потребуется?
    \item Сколько в среднем бросков ему потребовалось, если ни разу не выпала нечётная грань?
    \item Сколько в среднем бросков ему потребовалось, если нечётная грань выпала $Y$ раз?
\end{enumerate}

\item Альберт Ширяев подкидывает кубик до первой шестёрки. 
За каждый бросок казино платит Альберту один рубль. 
Однако, если нечётная грань выпадет хотя бы один раз, то Альберт не получает ничего. 
\begin{enumerate}
    \item Какова вероятность того, что Альберт получит положительный выигрыш?
    \item Сколько в среднем денег заработает Альберт?
\end{enumerate}

\end{enumerate}


\end{document}

